\section{Background and related work}\label{sec:background}

The proposed Self-Aware Steel Structure builds on three mature bodies of knowledge that have so far evolved largely in parallel. Structural health monitoring (SHM) is the discipline for inferring structural condition from measured response. Data-driven modeling compresses high-dimensional measurements into interpretable low-dimensional representations of system behavior. Machine learning (ML) supplies a practical toolset for civil and structural engineering. This section reviews these foundations and clarifies the terminology and assumptions on which the remainder of the paper relies. The recurring theme is that each field provides a necessary but individually insufficient ingredient: SHM defines \textit{what} must be inferred and \textit{why} a baseline comparison is required, data-driven modeling defines \textit{how} raw multi-channel signals are reduced to features that a digital representation can reason about, and ML defines \textit{how} those features are mapped to actionable statements about structural integrity \citep{farrar2013shm,brunton2019data,naser2023ml}. The novelty of the present work, developed in later sections, lies in fusing these ingredients into a structure that interprets its own state rather than merely transmitting data.

\subsection{Structural health monitoring as statistical pattern recognition}\label{subsec:shm-spr}

The canonical formulation of SHM frames damage detection as a statistical pattern recognition (SPR) problem \citep{farrar2013shm}. In this view, SHM is the process of implementing a damage detection strategy by observing a structure over time through periodically spaced response measurements, extracting damage-sensitive features from those measurements, and statistically analyzing the features to determine the current state of system health \citep{farrar2007introduction}. The process is organized into four integrated parts: (i) operational evaluation, which establishes the life-safety and economic justification for monitoring, defines what constitutes damage for the system at hand, and identifies the operational and environmental conditions under which data will be acquired; (ii) data acquisition, which selects excitation methods, sensor types, numbers and locations, and the hardware for acquisition, storage and transmittal; (iii) feature selection and extraction, which identifies low-dimensional quantities that are sensitive to damage yet robust to confounding influences; and (iv) statistical modeling for feature discrimination, in which learning algorithms operate on the extracted features to quantify the damage state \citep{farrar2013shm}. Data normalization, cleansing, compression and fusion are embedded within these parts, and ML enters primarily in the third and fourth.

Two foundational principles underpin this paradigm and are inherited directly by the present framework. First, damage assessment is intrinsically a \textit{comparison between two states} of the system, one of which is assumed to represent an undamaged baseline; damage itself is defined as a change to the material or geometric properties of a system, including changes to boundary conditions and connectivity, that adversely affects current or future performance \citep{farrar2013shm}. The loosening of a bolted connection is an explicit example: it alters connectivity and typically adds energy dissipation, which manifests as an increase in measured damping. Second, sensors cannot measure damage directly. Feature extraction and statistical classification are required to convert sensor data into damage information, and without intelligent feature extraction the most damage-sensitive measurements are also the most sensitive to benign operational and environmental variability \citep{farrar2013shm}.

The depth of diagnosis achievable is organized by the Rytter hierarchy, which orders the damage-identification questions by increasing knowledge of the damage state: existence, location, type, severity (extent), and prognosis \citep{farrar2013shm}. This hierarchy maps onto the supervised/unsupervised distinction that governs the choice of learning algorithm. Unsupervised learning, which builds a model of a single healthy class and tests new data for consistency with that class, can typically establish the \textit{existence} of damage and sometimes its \textit{location}; this leads naturally to novelty and outlier detection, the appropriate paradigm when only undamaged data are available, as is usually the case because deliberately damaging high-value structures is impractical \citep{sohn2003statistical,wangcha2022unsupervised}. Determining the \textit{type} and \textit{severity} of damage generally requires supervised learning trained on labeled examples from both healthy and damaged states \citep{farrar2013shm,figueiredo2010machine}. A central practical challenge is data normalization: separating response changes caused by benign operational and environmental variability, such as temperature-induced stiffness change, from those caused by damage, so that environmental fluctuation is not mistaken for structural deterioration and reported as a false positive \citep{figueiredo2010machine}. The classification of healthy versus anomalous behavior must therefore minimize both false-positive and false-negative diagnoses, whose costs may be weighted very differently \citep{farrar2013shm}.

\subsection{Data-driven modeling and reduced-order representations}\label{subsec:datadriven}

Measurements from distributed sensor networks are high-dimensional, yet the dynamics they capture typically evolve on a low-dimensional manifold, exhibiting dominant low-dimensional patterns \citep{brunton2019data}. The singular value decomposition (SVD) provides a systematic, numerically stable, and entirely data-driven way to extract these dominant correlations, yielding a hierarchical representation in which the left singular vectors are the dominant spatial patterns and the ordered singular values quantify their energy content. Principal component analysis (PCA) and proper orthogonal decomposition (POD) are computed directly from the SVD of the mean-subtracted data matrix, and the SVD furnishes the optimal low-rank approximation of a data set \citep{brunton2019data}. For the present framework, this supplies the mechanism for compressing distributed strain and acceleration streams into a few dominant modes and for defining a compact, low-dimensional baseline of healthy behavior against which deviations can be flagged.

Several complementary data-driven methods extend this idea toward dynamics. Dynamic mode decomposition (DMD) identifies spatio-temporal coherent structures from measurement snapshots and, unlike SVD/POD, associates each mode with a temporal eigenvalue encoding an oscillation frequency and a growth or decay rate; DMD is purely data-driven, requires no knowledge of the governing equations, and applies equally to experimental and numerical data \citep{brunton2019data}. The sparse identification of nonlinear dynamics (SINDy) complements this by discovering interpretable governing equations through sparse regression over a candidate function library, recovering the few active terms that describe the dynamics \citep{brunton2019data}. Reduced-order models (ROMs) project high-dimensional dynamics onto a few dominant modes for rapid evaluation, and can be coupled with sparse, optimally placed sensors to reconstruct the full field from a limited number of measurements, which provides a principled basis for numerical-analysis-driven sensor placement \citep{brunton2019data,tan2020sensorplacement,hassani2023sensorplacement}. Autoencoders learn a nonlinear latent representation by mapping an input through an encoder to a latent space and reconstructing it through a decoder; a linear autoencoder reduces to the SVD/POD basis, connecting deep learning to PCA, and the reconstruction error of a model trained on healthy data provides a natural residual that grows under abnormal conditions, an indicator well suited to anomaly detection \citep{brunton2019data}. Recurrent architectures, in particular long short-term memory (LSTM) networks, are designed to exploit the temporal structure of sequential data and can learn and reproduce the evolution of dynamical systems from short windows of observations \citep{brunton2019data}.

\subsection{Machine learning in civil and structural engineering}\label{subsec:ml-civil}

Within civil engineering, ML is best understood as a staged workflow comprising data collection, model development (training), model evaluation, and model optimization, in which an algorithm learns the mapping between input features and a target response \citep{naser2023ml}. This framing explicitly motivates SHM applications, in which sensor networks continuously measure and record information over the service life of a structure. Problems are partitioned into supervised learning, encompassing regression for numeric targets such as stress or displacement and classification for categorical targets such as discrete damage levels (for example, minor damage, major damage, and collapse), and unsupervised learning, which discovers the underlying structure of unlabeled data through clustering, dimensionality reduction, and anomaly detection \citep{naser2023ml}. Because labeled damaged-state data are scarce in SHM, semi-supervised variants and synthetic data generation are relevant, as is the use of a validated finite-element model run as a parametric study to generate augmented, hybrid (real plus simulated) datasets for training, although experimental data are generally favored over purely synthetic data \citep{naser2023ml,karaaslan2021attention}.

The algorithms named in the proposed AI layer are well established in this context. Random Forest is a non-parametric ensemble of decision trees that aggregates outcomes by averaging for regression or majority voting for classification and makes no assumption about the functional form of the feature--target relationship \citep{naser2023ml}. Support vector machines separate classes using a margin-maximizing hyperplane and accommodate nonlinearity through kernel functions such as the radial basis function, with a regression variant (SVR) \citep{naser2023ml}. Artificial neural networks arrange neurons in layers connected through differentiable activation functions and are trained by back-propagation and gradient descent; networks with multiple hidden layers constitute deep learning, and recurrent networks address sequential data, with LSTM motivated by the vanishing-gradient problem of plain recurrent networks \citep{naser2023ml}. Rigorous evaluation is essential: reporting several independent performance and error metrics, using confusion matrices, ROC curves, precision, recall, F1 and balanced accuracy for classification, and adopting stratified $k$-fold cross-validation guards against the bias of simple ratio sampling and against overfitting and underfitting \citep{naser2023ml}. The trustworthiness of an autonomous integrity assessment depends on explainability: permutation-based feature importance, partial dependence plots, SHAP and LIME provide traceable, physics-checkable explanations of why a model reaches a given diagnosis, moving beyond black-box prediction toward interpretable damage indicators \citep{naser2023ml}. Together, these methods supply the analytical machinery on which the cognitive layer of the proposed framework is built, while the data-scarcity remedies and explainability tools directly address the practical obstacles to deploying it on real steel infrastructure \citep{avci2021review}.

\subsection{Structural dynamics and vibration-based monitoring}\label{subsec:dynamics}

The premise that links a steel structure's measurable vibration response to its mechanical condition rests on the classical theory of structural dynamics. For a single-degree-of-freedom (SDOF) idealization, the equation of motion assembles inertia, damping and elastic restoring forces under dynamic equilibrium, $m\ddot{u} + c\dot{u} + ku = p(t)$ \citep{chopra2017dynamics}. Free undamped vibration of this system is simple harmonic motion whose natural circular frequency is $\omega_n = \sqrt{k/m}$, with the corresponding natural period $T_n = 2\pi/\omega_n$ and cyclic frequency $f_n = \omega_n/2\pi$. A property of central importance to monitoring follows directly: the natural vibration characteristics depend only on the mass and stiffness of the structure and not on the excitation. Consequently, of two systems of equal mass the stiffer one exhibits the higher natural frequency, and any reduction in stiffness, whether from yielding, cracking, support degradation or bolted-connection loosening, is in principle observable as a downward shift in the measured natural frequencies \citep{chopra2017dynamics,avci2021review}.

Real steel structures possess many degrees of freedom, and their dynamics are described by the matrix eigenvalue problem $\mathbf{k}\boldsymbol{\phi}_n = \omega_n^2 \mathbf{m}\boldsymbol{\phi}_n$, whose nontrivial solutions require $\det(\mathbf{k} - \omega_n^2\mathbf{m}) = 0$ \citep{chopra2017dynamics}. Because the mass and stiffness matrices are symmetric and positive definite, this characteristic equation yields $N$ real positive eigenvalues, ordered $\omega_1 < \omega_2 < \cdots < \omega_N$, each associated with a natural mode shape $\boldsymbol{\phi}_n$ describing a fixed deflected configuration in which all degrees of freedom vibrate in phase. The modes are orthogonal with respect to both mass and stiffness, which permits the coupled equations of a classically (proportionally) damped system to be transformed into $N$ independent modal SDOF equations and the total response to be recovered by modal superposition \citep{chopra2017dynamics}. The compact triplet of modal parameters, natural frequencies $\omega_n$, mode shapes $\boldsymbol{\phi}_n$ and modal damping ratios $\zeta_n$, therefore constitutes a physically interpretable fingerprint of the structure's current stiffness and mass distribution, and it is precisely this fingerprint that a self-aware structure must continuously identify and compare against a healthy baseline.

Damping characterizes energy dissipation through the ratio $\zeta = c/c_{cr}$; for most real structures it lies below 20\% and typically amounts to only a few percent, slightly reducing the damped frequency to $\omega_D = \omega_n\sqrt{1-\zeta^2}$ \citep{chopra2017dynamics}. Damping is generally not deducible from geometry and material alone but determined experimentally, for instance from the half-power (3~dB) bandwidth of the frequency-response curve, $\zeta = (f_b - f_a)/2f_n$, which is evaluable without knowledge of the applied force magnitude. Measured damping is also amplitude-dependent and tends to increase when a structure is driven into damaging regimes, a behavior with direct consequences for distinguishing healthy from abnormal states. The empirical precedent of the Millikan Library is instructive: at the larger amplitudes of the San Fernando earthquake its fundamental periods lengthened by roughly 50\%, which Chopra explicitly interprets as a reduction in stiffness accompanied by a substantial increase in damping, with apparent stiffness recovery (shorter periods) afterward \citep{chopra2017dynamics}. This documents, in a textbook setting, that monitored modal parameters track changes in structural condition.

The diagnostic is operationalized through forced harmonic vibration. Under a harmonic force $p(t) = p_o\sin\omega t$, the steady-state amplitude is governed by the deformation response factor and a phase angle, and the largest response occurs at resonance, where the forcing frequency approaches $\omega_n$; for light damping the resonant amplification is approximately $1/2\zeta$, so even small damping yields large amplitudes whose sharpness encodes $\zeta$ \citep{chopra2017dynamics}. A rotating-unbalance device with eccentric mass provides exactly such controlled excitation, producing a force $p(t) = (m_e e\,\omega^2)\sin\omega t$ whose amplitude scales with $\omega^2$; sweeping the forcing frequency through $\omega_n$ traces the frequency-response curve, whose peak locates the natural frequency and whose evolution reveals stiffness- or damage-induced shifts. Figure~\ref{fig:frf} sketches this behavior, showing how the resonance peak migrates to a lower frequency as stiffness is lost. Experimental modal testing and system identification thus convert raw acceleration and strain signals into the modal feature set that anchors vibration-based monitoring \citep{yu2024steelbridge}.

\subsection{Finite element models, digital twins and Industry 4.0}\label{subsec:fem-dt}

Whereas modal theory specifies \textit{what} to monitor, the finite element method (FEM) supplies the numerical model against which measured behavior is compared. FEM approximates a continuous field in piecewise fashion over small subregions using interpolation functions; each element carries a characteristic stiffness matrix that, on assembly, forms the global stiffness matrix relating nodal degrees of freedom to nodal loads, $[\mathbf{K}]\{\mathbf{D}\} = \{\mathbf{R}\}$ \citep{cook2002fea}. Skeletal steel systems such as the braced frame of square hollow sections envisioned here are naturally idealized with bar and beam/frame elements: the two-node bar carries axial stiffness $AE/L$, while plane- and space-frame elements add flexural and torsional degrees of freedom through coordinate transformation from local to global axes. Representing bolted beam--column connections as pinned, rigid, or semi-rigid is a deliberate modeling decision, and the adequacy of boundary conditions must be verified, since insufficient restraints leave $[\mathbf{K}]$ singular or near-singular \citep{cook2002fea}.

For dynamics, FEM produces the discrete equation of motion $[\mathbf{M}]\{\ddot{\mathbf{D}}\} + [\mathbf{C}]\{\dot{\mathbf{D}}\} + [\mathbf{K}]\{\mathbf{D}\} = \{\mathbf{R}\}$, with mass represented by consistent or lumped matrices and damping commonly modeled as Rayleigh (proportional) damping, $[\mathbf{C}] = a[\mathbf{M}] + b[\mathbf{K}]$, which preserves modal orthogonality \citep{cook2002fea}. Omitting damping yields the undamped free-vibration eigenproblem $([\mathbf{K}] - \omega^2[\mathbf{M}])\{\mathbf{D}\} = \{\mathbf{0}\}$, in which $[\mathbf{K}] - \omega^2[\mathbf{M}]$ is the dynamic stiffness matrix and a mode is a configuration where elastic resistance balances inertia. The resulting natural frequencies and mode shapes form the baseline for harmonic (frequency-response) and transient analyses; the textbook explicitly identifies rotating machinery attached to a structure as a source of harmonic excitation, mapping directly onto the planned testbed, and notes that plotting forced-vibration amplitude against excitation frequency reveals resonance peaks near the natural frequencies \citep{cook2002fea}. Because eigenvalue and eigenvector changes before and after a structural modification can be assessed without full re-solution, the FEM provides a tractable mechanism for relating observed modal shifts to localized stiffness loss.

Critically for any digital-twin claim, Cook et al. draw a sharp epistemic boundary: ``FEA is simulation, not reality'' \citep{cook2002fea}. The analysis is performed on a mathematical idealization, and even highly accurate computation can disagree with the physical structure when the model is inadequate. They classify the error sources as modeling error, discretization error and numerical error, advocate convergence checks through mesh refinement, and recommend obtaining an independent solution (analytical, handbook or experimental) to verify computed results. This same logic underwrites model updating: experiment is a trusted means of checking and improving the mathematical model, and reduced-order representations via static condensation and substructuring make repeated, near-real-time evaluation feasible. The systematic correlation of FE predictions with measured response, and the flagging of deviations as indicators of model inadequacy, is precisely the conceptual basis on which a digital twin operates.

A digital twin extends this calibrated FEM into a living, continuously updated counterpart of the physical asset, synchronized by streaming sensor data so that experimental responses are compared in real time against numerical predictions and deviations are detected through statistical indicators \citep{chiachio2022structural,mariniello2025enhancing}. Positioned within the broader Industry 4.0 paradigm of cyber-physical systems, the Industrial Internet of Things, pervasive embedded sensing and edge/cloud analytics, such a twin transforms the FEM from a one-off design artifact into a persistent, evolving service-life model of the structure \citep{zou2026deep,li2025edge,he2022integrated}. In this setting the validated FEM is not merely a predictive tool but the physics-bearing anchor that gives meaning to sensor streams, supplying the baseline modal and stress fields against which a self-aware steel structure can continuously interpret its own state, the integration of which with the data-driven and machine-learning layers is developed in the sections that follow.

\subsection{Positioning against existing digital-twin and machine-learning frameworks}\label{subsec:positioning}

The ingredients assembled above have already been combined in several recent monitoring frameworks, and the proposal of this paper must be read against them rather than in isolation. Chiachio et al. formalise a structural digital twin and its technology integration, but treat interpretation as a model-updating and prediction service rather than an embedded, autonomous capability of the structure \citep{chiachio2022structural}. Zou et al. couple deep learning with a digital twin for masonry SHM and candidly frame the open challenges, yet target neither steel framing nor a graded notion of autonomy \citep{zou2026deep}. Li et al. push autonomy furthest with an edge--fog--cloud architecture for distributed monitoring of a dam cluster, although autonomy there is a property of the computing fabric rather than of the structure's self-interpretation \citep{li2025edge}. Mariniello et al. add data-integrity governance to a digital-twin pipeline through blockchain and smart contracts \citep{mariniello2025enhancing}, which is complementary to, but distinct from, governing how a healthy baseline is allowed to evolve. On the algorithmic axis, Karyofyllas et al. propose a digital-twin-driven machine-learning framework for condition monitoring \citep{karyofyllas2025dtml}, Zhang and Sun integrate physics-guided learning with finite-element model updating \citep{zhang2021physicsguided}, and Parisi et al. localise damage on a steel structure directly from raw strain data \citep{parisi2022steel}; each strengthens one axis, model fusion, physics guidance, or steel-specific evidence, without unifying them under a single operational notion of structural self-awareness.

Read together, and as summarised in Table~\ref{tab:positioning}, these works show that the individual capabilities are mature, while three elements remain absent from any single framework: an explicit, graded scale of how autonomously a structure interprets its own state, the L0--L5 maturity model of Section~\ref{sec:paradigm}; a governed feedback policy that lets the healthy baseline adapt to benign change while refusing to absorb slow degradation, developed in Section~\ref{subsec:decision}; and a steel-specific demonstrator and protocol that bind these to a concrete class of structures. The contribution of this paper is therefore not the integration of digital twins, finite-element models and machine learning, which is already established, but the organisation of that integration around structural cognition, through a maturity scale, a governance policy and a steel-oriented protocol that together turn a richly instrumented structure into one that can interpret and report its own condition.
